\documentclass[10pt,a4paper]{amsart}
\usepackage[margin=19mm]{geometry}
\usepackage{amsmath,amssymb,mathtools}
\usepackage[scr=rsfs,cal=euler]{mathalfa}
\usepackage{xfrac}
\usepackage[unicode,hidelinks,bookmarks=false]{hyperref}
\newcommand{\ie}{i.e.\ }
\newcommand{\cf}{cf.\ }
\newcommand{\resp}{respectively\ }
\newcommand{\Subsection}{\subsection}
\providecommand{\nobreakdash}{\nobreak\hbox{-}}
\setlength{\emergencystretch}{2em}
\multlinegap0pt
\usepackage[shortlabels]{enumitem}
\usepackage{amssymb}
\usepackage{tcolorbox}
\newtheorem{lemma}{Lemma}
\newcommand{\C}{\mathbb{C}}
\renewcommand{\a}{\alpha}
\newcommand{\id}{\text{id}}
\renewcommand{\v}{\mathbf{v}}
\title{A relationship between the Kauffman bracket skein algebras and Roger-Yang skein algebras of some small surfaces}
\author{Chloe Marple, Helen Wong}
\date{Source: arXiv:2510.23865v2 (CC BY 4.0)}
\begin{document}
\maketitle

\noindent\emph{Source and licence.} A relationship between the Kauffman bracket skein algebras and Roger-Yang skein algebras of some small surfaces. Version arXiv:2510.23865v2 (CC BY 4.0), distributed by arXiv under the Creative Commons Attribution 4.0 International licence, http://creativecommons.org/licenses/by/4.0/, as recorded in the arXiv licence metadata for this exact version and shown on the abstract page https://arxiv.org/abs/2510.23865v2. Full licence notice: \texttt{licenses/arxiv-2510.23865-CC-BY-4.0.txt}.\par\smallskip

\noindent\textit{Editorial scope.} Counted unit: the article's lemma lem:calculate-u (source0.tex lines 703--705). The article's complete original proof of it is delivered with it (source0.tex lines 706--762, one \texttt{proof} environment of 57 lines). No mathematics is written, repaired or re-derived: the statement and the proof are the article's own lines, in the article's own notation, and no retained line is substituted or re-flowed.  Retained uncounted as the source-local context the proof needs: lines 677--679.  Nothing was omitted from any retained span, so the package carries no omission record.  No retained line contains a citation, so the wrapper carries no bibliography and no bibliography entry.  No retained line contains a figure environment, an image-inclusion command or drawing code, so no image is delivered and the package declares no asset dependency.  Editorial formatting only, in addition: an AMS \texttt{amsart} preamble that restates the article's own declarations and macros used by the retained lines, namely the environments \texttt{lemma} on separate counters, together with the standard packages those lines need.  The retained environments print in this document as Lemma 1 in order of appearance, whereas the article prints the same items with the numbering of its own full document.  The complete native source of the article as distributed in the arXiv e-print for this exact version is bundled unchanged as \texttt{sources/187912/source0.tex}.
\begin{lemma}\label{lem:eigenspace}
The space $V$ is $N$ dimensional and admits a basis $\{\v_1,\ldots,\v_N\}$ where each $\v_k$ generates the eigenspace $V_k$.
\end{lemma}

\begin{lemma}\label{lem:calculate-u}
	We have $E=(v_1v_2)^{-N}(t_1^2 +t_2^2 +t_1t_2t_3)$ and $u=-\dfrac{t_1+x^Nt_2}{\sqrt{v_1v_2}^N}.$
\end{lemma}

\begin{proof}
    Define $U,D:W\to W$ for $\v_k\in V_k$ as
\[        U\v_k=-\frac{x^{-1}A^{-2k-1}}{xA^{2k}-x^{-1}A^{-2k}}U_k\v_k:=H^U_k\cdot U_k\v_k
\qquad \text{and}\qquad
        D\v_k=\frac{xA^{2k-1}}{xA^{2k}-x^{-1}A^{-2k}}D_k\v_k:=H_k^D\cdot D_k\v_k
 \]
    and notice $U(V_k)=V_{k+1}$ and $D(V_k)=V_{k-1}$. Also, note $\rho(\beta)=U+D$. Hence,
    \[
    t_1\id_V=T_N(\rho(\beta)\sqrt{v_1v_2})=T_N(\sqrt{v_1v_2}(U+D))
    \]
    The left hand side is a polynomial in $\sqrt{v_1v_2}U$ and $\sqrt{v_1v_2}D$ over $\C[v_1,v_2]$. Consider a monomial of length $m$, where $n$ of the terms are $\sqrt{v_1v_2}U$ and $m-n$ are $\sqrt{v_1v_2}D$. It is a property of $T_N$ that all the monomials are of odd degree, so $m$ is odd. The monomial sends $V_k$ to $V_{k+n-(m-n)}=V_{k+2n-m}$. As $T_N(\sqrt{v_1v_2}(U+D))=t_1\id_V$ is homothety, it must fix all the $V_k$. Hence, the only monomials with nonzero coefficient must satisfy $2n-m\equiv 0\text{ mod }N$. As $n,m\in \{0,\ldots,N\}$, we have $(m,n)$ is $(N,0)$ or $(N,N)$. As $T_N(\sqrt{v_1v_2}(U+D))$ has degree $N$, $T_N(\sqrt{v_1v_2}(U+D))=\sqrt{v_1v_2}^{N}U^N+\sqrt{v_1v_2}^{N}D^N$. 
   
    Let $1\leq k\leq N$. One can calculate
    \begin{align*}
        U^N\v_k&=U^{N-1}(H_k^UU_k\v_k)&\\
        &=H_kH^U_{k+1}\cdots H^U_{k+N-1}U_{k+N-1}\cdots U_k\v_k&\\
        &=H_kH^U_{k+1}\cdots H^U_{k+N-1}u\v_k &\text{from the formula for $U_k\v_k$ in Lemma \ref{lem:eigenspace}}
    \end{align*}
    Also, as $N$ is odd, and $A^2$ is a primitive $N$th root of unity, we have $\prod_{k=1}^N(xA^{2k}-x^{-1}A^{-2k})=x^N-x^{-N}$. Hence,
    \begin{align*}
        H^U_kH^U_{k+1}\cdots H^U_{k+N-1}&=\prod_{k=1}^NH^U_k
        =\prod_{k=1}^N(-\frac{x^{-1}A^{-2k-1}}{xA^{2k}-x^{-1}A^{-2k}})
        =-\frac{x^{-N}A^{-N2k}A^{-N}}{\prod_{k=1}^N(xA^{2k}-x^{-1}A^{-2k})}
        =\frac{x^{-N}}{x^N-x^{-N}}
    \end{align*}
    Thus, $U^N\v_k=\frac{x^{-N}u}{x^N-x^{-N}}\v_k.$
    We can also calculate $D^N\v_k=H_k^DH_{k+1}^D\cdots H_{k-N+1}^DD_{k-N+2}\cdots D_k\v_k$.
    As before,
    \[
    H_k^DH_{k+1}^D\cdots H_{k+N-1}^D=\frac{x^NA^{2kN}A^{-N}}{\prod_{k=1}^N (xA^{2k}-x^{-1}A^{-2k})}
    =-\frac{x^N}{x^N-x^{-N}}
    \]
    and, by applying the formula for $D_k\v_k$ in Lemma \ref{lem:eigenspace},
    \begin{align*}
        \underbrace{D_{k-N+1}D_{k-N+2}\cdots D_{k-1}D_k}_{N \text{ terms}}\v_k
        =u^{-1}E\v_k
    \end{align*}
    Hence, $D^N\v_k=-\frac{u^{-1}x^NE}{x^N-x^{-N}}\v_k$.
    So we have
    \begin{equation}\label{eqn:t1-u-poly}
     t_1\v_k=T_N(\sqrt{v_1v_2}(U+D))\v_k
    =(\sqrt{v_1v_2}^{N}U^N+\sqrt{v_1v_2}^{N}D^N)\v_k
    =\sqrt{v_1v_2}^{N}\frac{x^{-N}u-x^Nu^{-1}E}{x^N-x^{-N}}\v_k
    \end{equation}
    which determines $u$ up to two possibilities.

    To determine $u$ completely, we do the same calculations with $\rho(\a)$. Observe $\rho(\a)=xA^{2k+1}U+x^{-1}A^{1-2k}D$. As before,
    \[
    t_2\id_V=T_N(\rho(\a)\sqrt{v_1v_2})=T_N(\sqrt{v_1v_2}(xA^{2k+1}U+x^{-1}A^{1-2k}D))=-\sqrt{v_1v_2}^{N}x^NU^N-\sqrt{v_1v_2}^{N}x^{-N}D^N.
    \]
    So by the calculations above of $U^N$ and $D^N$,
    \begin{equation}\label{eqn:t2-u-poly}
    t_2=\sqrt{v_1v_2}^{N}\frac{-u+u^{-1}E}{x^N-x^{-N}}.
    \end{equation}
    
    Now, $u$ is a common root of \eqref{eqn:t1-u-poly} and \eqref{eqn:t2-u-poly}. By explicitly finding the roots of these equations and a bit of casework, one can show that they have common roots if and only if $E=(v_1v_2)^{-N}(t_1^2+t_2^2+t_1t_2t_3)$, and, in this case, $-\frac{t_1+x^Nt_2}{\sqrt{v_1v_2}^{N}}$  is always a common root, and it is unique when $u\neq0$, which is the case here.
\end{proof}

\end{document}
